% Part: many-valued-logic
% Chapter: three-valued-logics
% Section: lukasiewicz

\documentclass[../../../include/open-logic-section]{subfiles}

\begin{document}

\olfileid{mvl}{thr}{luk}

\olsection{Łukasiewicz逻辑（卢卡西维茨逻辑）}

最早发表且成体系的多值逻辑提案之一，要归功于波兰哲学家 Jan \L ukasiewicz 在 1921 年的工作。\L ukasiewicz 的动机来自亚里士多德的海战问题：似乎在\emph{今天}，“明天将会有一场海战”这个语句既不真也不假：它的真值尚未确定。\L ukasiewicz 提议对这种“未来偶然”语句引入第三个真值。

\begin{quote}
我可以毫无矛盾地假设，我在明年某个特定时刻（例如 12 月 21 日正午）是否在华沙，在目前既没有肯定地也没有否定地确定。因此，我在那个特定时间出现在华沙是可能的，但并非必然的。在这一假设下，“我将于明年 12 月 21 日正午在华沙”这个命题在目前时刻既不真也不假。因为如果它现在为真，那么我将来到华沙的未来事件就是必然的，这与假设相矛盾。另一方面，如果它现在为假，那么我将来到华沙的未来事件就是不可能的，这也与假设相矛盾。因此，所考虑的命题在当下既不真也不假，并且必定拥有一个不同于“0”（假）和“1”（真）的第三值。我们可以用“$\frac{1}{2}$”来指代这个值。它代表“可能”，并与“真”和“假”一起作为第三个真值。
\end{quote}

我们将用 $\Undef$ 表示 \L ukasiewicz 的第三个真值。\footnote{\L ukasiewicz 在此使用“可能”的方式在今天并不常见，即用它来表示“可能但非必然”。}

在这种解释下，联结词 $\lnot$、$\land$ 和 $\lor$ 的真值函数不难确定：一个未来偶然语句的否定也是一个未来偶然语句，所以 $\tf{\lnot}(\Undef) = \Undef$。如果一个合取式的一个合取支是未确定的而另一个为真，那么该合取式也是未确定的——毕竟，取决于那个未来偶然的合取支最终如何，这个合取式可能为真，也可能为假。所以 \[
    \tf{\land}(\True, \Undef) =
\tf{\land}(\Undef, \True) =
\Undef.
\]
然而，如果另一合取支为假，那么该合取式就不可能为真，所以 \[\tf{\land}(\False, \Undef) =
\tf{\land}(\False, \Undef) = \False.\]
其他取值情况（若自变元取已确定的真值，即 $\True$ 或 $\False$）则与经典逻辑相同。

对于条件句，情况就有点棘手了。假设 $q$ 是一个未来偶然语句。如果 $p$ 为假，那么无论 $q$ 的结果如何，$p \lif q$ 都将为真，因此我们应设 $\tf{\lif}(\False, \Undef) = \True$。而如果 $p$ 为真，那么无论 $q$ 的结果如何，$q \lif p$ 都将为真，所以 $\tf{\lif}(\Undef, \True) = \True$。如果 $p$ 为真，那么 $p \lif q$ 可能为真也可能为假，因此 $\tf{\lif}(\True,
\Undef) = \Undef$。类似地，如果 $p$ 为假，那么 $q \lif p$ 可能为真也可能为假，因此 $\tf{\lif}(\Undef, \False) = \Undef$。剩下的情况是 $p$ 和 $q$ 都是未来偶然语句。基于前述动机，我们确实应该在这种情况下赋值为 $\Undef$。然而，这会导致 $!A \lif !A$ 不再是一个重言式。\L ukasiewicz 并不介意放弃 $!A \lor \lnot !A$ 和 $\lnot(!A \land \lnot !A)$，但他不愿放弃 $!A \lif !A$。因此他规定 $\tf{\lif}(\Undef, \Undef) =
\True$。

\begin{defn}\ollabel{def:lukasiewicz}
三值 \L ukasiewicz 逻辑由如下矩阵定义：
\begin{enumerate}
  \item 使用标准命题语言 $\Lang L_0$，包含联结词 $\lnot$、$\land$、$\lor$、$\lif$。
  \item 真值集合为 $V = \{\True, \Undef, \False\}$。
  \item $\True$ 是唯一的指定值，即 $V^+ = \{\True\}$。
  \item 真值函数由下列表格给出：
  \begin{center}
    \begin{tabular}{c|c}
      $\tf{\lnot}$ & \\
      \hline
      $\True$ & $\False$ \\
      $\Undef$ & $\Undef$ \\
      $\False$ & $\True$
    \end{tabular}
    \quad
    \begin{tabular}{c|ccc}
      $\tf{\land}[\LogLuk[3]]$ & $\True$ & $\Undef$ & $\False$ \\
      \hline
      $\True$ & $\True$ & $\Undef$ & $\False$ \\
      $\Undef$ & $\Undef$ & $\Undef$ & $\False$\\
      $\False$ & $\False$ & $\False$ & $\False$
    \end{tabular}
    \\[2ex]
    \begin{tabular}{c|ccc}
      $\tf{\lor}[\LogLuk[3]]$ & $\True$ & $\Undef$ & $\False$ \\
      \hline
      $\True$ & $\True$ & $\True$ & $\True$ \\
      $\Undef$ & $\True$ & $\Undef$ & $\Undef$ \\
      $\False$ & $\True$ & $\Undef$ & $\False$
    \end{tabular}
    \quad
    \begin{tabular}{c|ccc}
      $\tf{\lif}[\LogLuk[3]]$ & $\True$ & $\Undef$ & $\False$ \\
      \hline
      $\True$ & $\True$ & $\Undef$ & $\False$ \\
      $\Undef$ & $\True$ & $\True$ & $\Undef$  \\
      $\False$ & $\True$ & $\True$ & $\True$
    \end{tabular}
  \end{center}
\end{enumerate}
\end{defn}

可以轻易看出，任何仅包含 $\lnot$、$\land$ 和 $\lor$ 的公式~$!A$，如果它的所有命题变元都被赋值为~$\Undef$，那么该公式的真值也将是~$\Undef$。
因此，例如经典的重言式 $p \lor \lnot p$ 和 $\lnot(p \land \lnot p)$ 在 $\LogLuk[3]$ 中不再是重言式，因为只要 $\pAssign v(p) = \Undef$，就有 $\pValue{v}(!A) = \Undef$。

在 $\pAssign v(p)$ 为 $\True$ 或 $\False$ 的赋值下，$\pValue v(!A)$ 将与它的经典真值一致。

\begin{prop}
  如果 $\pAssign v(p) \in \{\True, \False\}$ 对~$!A$ 中所有的 $p$ 成立，那么
  $\pValue v(!A)[\LogLuk[3]] = \pValue v(!A)[\LogCL]$。
\end{prop}

\begin{prob}\label{mvl:thr:luk:prob:luk-iff} 假设我们在~$\LogLuk[3]$ 中定义 $\pValue v(!A \liff !B) = \pValue v((!A \lif !B) \land (!B \lif
  !A))$。那么 $\liff$ 的真值表会是怎样的？
\end{prob}

许多经典重言式在 \LogLuk[3] 中\emph{也}是重言式，例如 $\lnot p \lif (p \lif q)$。与经典逻辑一样，我们可以用真值表来验证这一点：

\begin{center}
\begin{tabular}{cc|cccccc}
  $p$ & $q$ & $\lnot$ & $p$ & $\lif$ & $\smash{(}p$ & $\lif$ & $q\smash{)}$ \\ \hline
  \True & \True & \False & \True & \True & \True & \True & \True \\
  \True & \Undef & \False & \True & \True & \True & \Undef & \Undef \\
  \True & \False & \False & \True & \True & \True & \False & \False \\
  \Undef & \True & \Undef & \Undef & \True & \Undef & \True & \True \\
  \Undef & \Undef & \Undef & \Undef & \True & \Undef & \True & \Undef \\
  \Undef & \False & \Undef & \Undef & \True & \Undef & \Undef & \False \\
  \False & \True & \True & \False & \True & \False & \True & \True \\
  \False & \Undef & \True & \False & \True & \False & \True & \Undef \\
  \False & \False & \True & \False & \True & \False & \True & \False \\
\end{tabular}
\end{center}

\begin{prob}
  证明下列公式在 \LogLuk[3] 中是重言式：
  \begin{enumerate}
    \item $p \lif (q \lif p)$
    \item\label{mvl:thr:luk:prob:luk-taut-2} $\lnot(p \land q) \liff (\lnot p \lor \lnot q)$
    \item\label{mvl:thr:luk:prob:luk-taut-3} $\lnot(p \lor q) \liff (\lnot p \land \lnot q)$
  \end{enumerate}
  （对于 \olref[mvl][thr][luk]{prob:luk-taut-2} 和 \olref[mvl][thr][luk]{prob:luk-taut-3}，将 $!A \liff !B$ 视为 $(!A \lif !B) \land (!B \lif !A)$ 的缩写，或参考你对 \olref[mvl][thr][luk]{prob:luk-iff} 的解答。）
\end{prob}

\begin{prob}
  证明下列经典重言式在~\LogLuk[3] 中不是重言式：
  \begin{enumerate}
    \item $(\lnot p \land p) \lif q)$
    \item $((p \lif q) \lif p) \lif p$
    \item $(p \lif (p \lif q)) \lif (p \lif q)$
  \end{enumerate}
\end{prob}

人们因此也许会认为，尽管并非所有经典重言式在~$\LogLuk[3]$ 中都是重言式，但它们至少在每个赋值下都应取真值~$\True$ 或~$\Undef$。事实并非如此。反例由 \[
  \lnot(p \lif \lnot p) \lor \lnot(\lnot p \lif p)
\] 给出，当 $p$ 为~$\Undef$ 时，该公式为 $\False$。

\begin{prob}
  下列关系在 \L ukasiewicz 逻辑中哪些成立？为每个关系给出一个真值表。
  \begin{enumerate}
    \item $p, p \lif q \Entails q$
    \item $\lnot\lnot p \Entails p$
    \item $p \land q \Entails p$
    \item $p \Entails p \land p$
    \item $p \Entails p \lor q$
  \end{enumerate}
\end{prob}

\L ukasiewicz 希望在他的三值系统基础上建立一种关于可能性的逻辑，方法是引入一元联结词 $\Diamond !A$（表示“$!A$ 是可能的”）以及相应的 $\Box !A$（表示“$!A$ 是必然的”）：

\begin{center}
  \begin{tabular}{c|c}
    $\tf{\Diamond}$ & \\
    \hline
    $\True$ & $\True$ \\
    $\Undef$ & $\True$ \\
    $\False$ & $\False$
  \end{tabular}
  \quad
  \begin{tabular}{c|c}
    $\tf{\Box}$ & \\
    \hline
    $\True$ & $\True$ \\
    $\Undef$ & $\False$ \\
    $\False$ & $\False$
  \end{tabular}
\end{center}

换言之，$p$ 是可能的，当且仅当它尚未被确定为假；
而 $p$ 是必然的，当且仅当它已被确定为真。

\begin{prob}
  验证 $\Box p \liff \lnot\Diamond \lnot p$ 和 $\Diamond p \liff
  \lnot \Box \lnot p$ 在扩展了上述 $\Box$ 与~$\Diamond$ 真值表的~\LogLuk[3]~中是重言式。
\end{prob}

然而，这种模态逻辑方案的缺陷很快就显现出来了：无论事情如何发展，$p \land \lnot p$ 都不可能为真。
所以，即使它现在未被确定（因而是未确定的），它也应该被视为不可能的，也就是说，$\lnot \Diamond(p
\land \lnot p)$ 应该是一个重言式。
但是，如果 $\pAssign v(p) = \Undef$，那么 $\pValue v(\lnot \Diamond(p \land \lnot p)) =
\Undef$。
尽管 \L ukasiewicz 认为两个真值不足以容纳可能性和必然性之类的模态区分这一点是正确的，但引入第三个真值依然不够。

\end{document}